\documentclass[conference]{IEEEtran}
\IEEEoverridecommandlockouts

\usepackage{cite}
\usepackage{amsmath,amssymb,amsfonts}
\usepackage{algorithmic}
\usepackage{graphicx}
\usepackage{textcomp}
\usepackage{xcolor}
\usepackage{booktabs}
\usepackage{array}
\usepackage{url}

\def\BibTeX{{\rm B\kern-.05em{\sc i\kern-.025em b}\kern-.08em
    T\kern-.1667em\lower.7ex\hbox{E}\kern-.125emX}}

\begin{document}

\title{Data-Driven Discovery of Economic Spatial Profiles\\
in Morocco Using Unsupervised Machine Learning}

\author{
\IEEEauthorblockN{1\textsuperscript{st} Souhaib Madhour}
\IEEEauthorblockA{\textit{Département Informatique} \\
\textit{EMSI Marrakech}\\
Marrakech, Morocco \\
souhaibmadhour@gmail.com}
\and
\IEEEauthorblockN{2\textsuperscript{nd} Safia Tirar}
\IEEEauthorblockA{\textit{Département Informatique} \\
\textit{EMSI Marrakech}\\
Marrakech, Morocco \\
safiatirar6@gmail.com}
\and
\IEEEauthorblockN{3\textsuperscript{rd} Yassine Nait Ider}
\IEEEauthorblockA{\textit{Département Informatique} \\
\textit{EMSI Marrakech}\\
Marrakech, Morocco \\
yassinay205@gmail.com}
\and
\IEEEauthorblockN{4\textsuperscript{th} Mohamed Erraji}
\IEEEauthorblockA{\textit{Département Informatique} \\
\textit{EMSI Marrakech}\\
Marrakech, Morocco \\
medxd2002@gmail.com}
}

\maketitle

% -------------------------------------------------------
\begin{abstract}
Understanding the economic structure of urban and peri-urban zones
is critical for regional planning, investment targeting, and
policy design in developing economies. In Morocco, no standardized,
fine-grained index of infra-regional economic attractiveness
currently exists. This paper proposes a data-driven pipeline
to discover, characterize, and assign economic spatial profiles
at the zone level using unsupervised and supervised machine learning.
Starting from a heterogeneous dataset of 37,873 formal business
entities geocoded across the Casablanca--Settat--El~Jadida region
and aggregated into a 1.5~km $\times$ 1.5~km spatial grid of 142 zones,
we apply a multi-algorithm clustering selection framework spanning
K-Means, Gaussian Mixture Models, Agglomerative Clustering, and
Spectral Clustering. The selected partition (Spectral Clustering,
$k=3$) achieves a balance score of 0.96 and reveals three
interpretable economic archetypes: \textit{High-density capital hubs},
\textit{Active emerging zones}, and \textit{Underdeveloped early-stage ecosystems}.
A Random~Forest classifier trained to predict discovered cluster labels
achieves a test macro-F1 of 0.955 and an accuracy of 0.966,
confirming the robustness and generalizability of the discovered profiles.
Feature importance analysis identifies business activity rate,
spatial density, entity counts, and sectoral diversification as
the primary structural drivers. The framework is reproducible,
interpretable, and extensible to other Moroccan regions,
offering a practical tool for data-informed territorial economic analysis.
\end{abstract}

\begin{IEEEkeywords}
economic zone profiling, unsupervised clustering, spatial machine learning,
Morocco, urban analytics, feature importance, spectral clustering,
territorial attractiveness
\end{IEEEkeywords}

% -------------------------------------------------------
\section{Introduction}

Territorial economic attractiveness has emerged as a central concern
for regional planners and policy-makers, particularly in rapidly
urbanizing economies of the Global South \cite{khalid2023artificial}.
Morocco, with its significant urban concentration and diverse regional
economic fabric, presents a compelling case: formal business registries
are increasingly digitized, yet no standardized infra-regional
economic profile framework exists at the zone level.

Traditional approaches to economic attractiveness rely on
aggregate administrative statistics (GDP per region, employment rates)
that mask substantial intra-regional heterogeneity
\cite{xu2023enhancing,feldmeyer2020using}.
Machine learning has been shown to uncover latent economic structures
from heterogeneous data sources in comparable urban analytics contexts
\cite{chen2025machine,huang2022machine,kourtit2021safe},
but applications targeting the Moroccan economic landscape
at high spatial resolution remain scarce.

This paper addresses the following scientific question:
\textit{Can data-driven unsupervised learning discover interpretable,
stable, and generalizable economic spatial profiles from raw business
registry data, without relying on a synthetic score as a training target?}

Our contributions are threefold:
\begin{itemize}
  \item A reproducible spatial aggregation pipeline converting
    individual business entities into zone-level economic feature vectors
    on a uniform 1.5~km $\times$ 1.5~km GPS grid.
  \item A multi-objective clustering selection framework comparing seven
    algorithm families across balance, separation, and supervisory readiness
    criteria, with Spectral Clustering selected as the optimal solution.
  \item A supervised profile assignment module achieving high macro-F1
    performance, enabling the assignment of new zones to discovered
    economic archetypes.
\end{itemize}

The remainder of this paper is organized as follows.
Section~\ref{sec:related} reviews related work.
Section~\ref{sec:data} describes the data collection and preprocessing.
Section~\ref{sec:method} presents the methodological pipeline.
Section~\ref{sec:results} reports experimental results and explainability analysis.
Section~\ref{sec:discussion} discusses implications and limitations.
Section~\ref{sec:conclusion} concludes the paper.

% -------------------------------------------------------
\section{Related Work}
\label{sec:related}

The intersection of machine learning and spatial economic analysis
has generated a growing body of work. Xu~et~al.~\cite{xu2023enhancing}
demonstrate that complex urban features extracted from heterogeneous
sources can significantly improve economic competitiveness analysis
using ML models. Similarly, Feldmeyer~et~al.~\cite{feldmeyer2020using}
show that OpenStreetMap data combined with ML can produce reliable
socio-economic indicators at fine spatial resolution, a finding
directly relevant to data-scarce national contexts.

In tourism and attractiveness estimation, Castellana~et~al.~\cite{castellana2024exploring}
apply ML to regional tourism attractiveness with counterfactual
explainability, while Giglio~et~al.~\cite{giglio2020machine}
combine points-of-interest data and ML to characterize typical
tourist Italian cities. Khalid~et~al.~\cite{khalid2023artificial}
specifically address small Moroccan cities, using AI for
attractiveness modeling and underscoring the feasibility of
ML-based approaches in the Moroccan administrative context.

Clustering-based spatial typology has been applied to urban vitality
\cite{huang2022machine}, safe city profiling \cite{kourtit2021safe},
and urban park attractiveness \cite{chen2025machine}.
These works consistently report that multi-variable clustering
over spatial economic features can reveal interpretable archetypes
that correspond to meaningful urban dynamics.

In contrast to supervised attractiveness prediction
(which faces the circular reasoning risk when the target is
constructed from the same input signals), our work adopts
an unsupervised-first framing, followed by supervised
profile assignment---a distinction explicitly recommended
in recent methodological literature \cite{agiropoulos2024exploring,sakti2024machine}.

% -------------------------------------------------------
\section{Data Collection and Preprocessing}
\label{sec:data}

\subsection{Data Sources}

Two primary sources were integrated. The Charika.ma business registry
provided 998,000 raw rows representing approximately 98,000 unique
formal entities, with multiple fiscal-year entries per company creating
a 90\% functional duplication rate. A two-pass deduplication was applied:
exact deduplication on (entity ID, trade name, city), followed by
approximate deduplication using Jaro-Winkler similarity with a threshold
of 0.92, chosen to balance over-merging (too low) and residual duplicates
(too high).

\subsection{Geographic Scope}

The operational perimeter was set to the Casablanca--Settat--El~Jadida
region, selected for its data concentration and economic coherence.
A triple geographic filter (GPS bounding box, normalized city name,
region code) reduced the initial 48,209 merged entities to 37,873
spatially valid entities, eliminating 10,336 out-of-region records.

\subsection{Spatial Aggregation}

Individual entities were aggregated into a uniform 1.5~km $\times$ 1.5~km
GPS grid, chosen to balance intra-urban resolution and minimum
statistical cell population. A critical preprocessing step addressed
a density artifact: 4,684 entities with missing GPS coordinates
had been imputed to the city centroid, creating an artificially dominant
cell. A binary flag \texttt{coords\_real} was introduced to distinguish
original from imputed coordinates; density features use only original
coordinates, while imputed entities contribute only to qualitative statistics.

After applying a minimum threshold of 3 real entities per cell,
the final dataset comprises \textbf{792 observations across 142 zones and 13 sectors},
forming the basis for all modeling experiments.

\subsection{Feature Engineering}

Twelve zone-level features were engineered from the aggregated data,
as summarized in Table~\ref{tab:features}.

\begin{table}[htbp]
\caption{Zone-Level Features Used in Clustering}
\label{tab:features}
\begin{center}
\begin{tabular}{|l|p{5cm}|}
\hline
\textbf{Feature} & \textbf{Description} \\
\hline
\texttt{entity\_count\_real} & Entities with real GPS coordinates \\
\texttt{entity\_count\_total} & All entities linked to the zone \\
\texttt{density\_log} & Log-transformed spatial density \\
\texttt{active\_rate} & Share of active entities \\
\texttt{capital\_median} & Median registered capital \\
\texttt{capital\_mean} & Mean registered capital \\
\texttt{capital\_max} & Maximum registered capital \\
\texttt{sector\_diversity} & Shannon-based sector diversification index \\
\texttt{formal\_ratio} & Formalization ratio (SARL/SA share) \\
\texttt{sarl\_count} & Number of SARL-type entities \\
\texttt{sa\_count} & Number of SA-type entities \\
\texttt{n\_sectors\_present} & Count of distinct active sectors \\
\hline
\end{tabular}
\end{center}
\end{table}

The synthetic attractiveness score, computed as a weighted composite
of these same features for descriptive validation purposes, was
\textit{explicitly excluded} from all clustering and classification
inputs to prevent circular reasoning---a key methodological
safeguard \cite{sakti2024machine}.

% -------------------------------------------------------
\section{Methodology}
\label{sec:method}

\subsection{Pipeline Architecture}

The modeling pipeline consists of four stages, as shown in Fig.~\ref{fig:pipeline}:
(1) spatial aggregation and feature construction; (2) multi-objective
clustering and algorithm selection; (3) cluster interpretation and
economic profiling; and (4) supervised profile assignment.

\begin{figure}[htbp]
\centerline{\includegraphics[width=\columnwidth]{artifacts/figures/morocco_cluster_map_typology.png}}
\caption{Spatial distribution of discovered economic profiles
  (typology mode, Spectral Clustering $k=3$) over the
  Casablanca--Settat--El~Jadida region. Each color denotes one
  economic archetype.}
\label{fig:pipeline}
\end{figure}

\subsection{Multi-Objective Clustering Selection}

Seven algorithm families were evaluated: K-Means ($k \in [3,6]$),
Gaussian Mixture Models with full and diagonal covariance
($k \in [3,6]$), Ward and average/complete Agglomerative Clustering
($k \in [3,6]$), and Spectral Clustering with nearest-neighbor
affinity ($k \in [3,6]$). DBSCAN was evaluated in a separate
anomaly mode described in Section~\ref{sec:anomaly}.

Algorithm selection was governed by a multi-objective score:
\begin{equation}
  \mathcal{O} = \alpha \cdot S + \beta \cdot B \cdot \mathbf{1}[\text{supervised-valid}]
  \label{eq:objective}
\end{equation}
where $S$ is the silhouette coefficient, $B$ is the entropy-based
balance score over cluster sizes, and $\mathbf{1}[\text{supervised-valid}]$
is an indicator requiring a minimum cluster share of 10\% for each cluster.
Weights $\alpha$ and $\beta$ were set to 0.3 and 0.7, prioritizing
a balanced and trainable partition over raw separation.

\subsection{Supervised Profile Assignment}

Once cluster labels were discovered, a supervised classifier was trained
to predict these labels on held-out zones, enabling generalization to
new geographic areas. Three classifiers were compared:
Random~Forest, XGBoost, and LightGBM, evaluated via 5-fold
cross-validated macro-F1. The classifier feature set extended
the clustering features with spatial coordinates (\texttt{cell\_lat},
\texttt{cell\_lon}) and a categorical city indicator.
A strict anti-leakage protocol excluded the attractiveness score
and cluster label from all predictor inputs.

\subsection{Anomaly Detection Mode}
\label{sec:anomaly}

In complementary anomaly mode, DBSCAN ($\varepsilon=2.0$,
\textit{min\_samples}=3) was applied to detect atypical
high-formalization or high-capital economic pockets.
This produced 6 noise-flagged zones and 3 compact clusters,
offering a secondary analytical lens on spatial economic outliers.

% -------------------------------------------------------
\section{Results and Explainability}
\label{sec:results}

\subsection{Clustering Selection}

Table~\ref{tab:clustering} summarizes the top-performing models
from the typology selection grid. Spectral Clustering with $k=3$
was selected with a silhouette of 0.229 and an exceptionally high
balance score of 0.963, satisfying the supervised readiness constraint.
K-Means yielded higher silhouette values but consistently produced
degenerate partitions with singleton clusters, failing the balance
and supervisory readiness criteria.

\begin{table}[htbp]
\caption{Clustering Algorithm Comparison (Top Results, Typology Mode)}
\label{tab:clustering}
\begin{center}
\begin{tabular}{|l|c|c|c|c|}
\hline
\textbf{Algorithm} & \textbf{$k$} & \textbf{Silhouette} & \textbf{Balance} & \textbf{Valid} \\
\hline
K-Means          & 3 & 0.266 & 0.665 & No \\
K-Means          & 6 & 0.353 & 0.598 & No \\
GMM (full)       & 3 & 0.277 & 0.700 & No \\
GMM (diag)       & 3 & 0.130 & 0.974 & \textbf{Yes} \\
Agglomerative (ward) & 6 & 0.330 & 0.734 & No \\
\textbf{Spectral}& \textbf{3} & \textbf{0.229} & \textbf{0.963} & \textbf{Yes} \\
Spectral         & 4 & 0.187 & 0.931 & Yes \\
\hline
\end{tabular}
\end{center}
\end{table}

\subsection{Discovered Economic Profiles}

The three discovered profiles exhibit coherent and distinguishable
economic characteristics:

\begin{itemize}
  \item \textbf{Cluster 0 --- Underdeveloped early-stage ecosystems}
    (29 zones, 20.4\%): Low density ($\log$-density $= 2.87$),
    low activity rate (0.58), and minimal formalization
    (formal ratio $= 0.017$). Post-hoc validation score: 35.0.
  \item \textbf{Cluster 1 --- Active emerging zones}
    (56 zones, 39.4\%): Moderate density ($\log$-density $= 2.42$),
    high activity rate (0.946), and rising formalization
    (formal ratio $= 0.173$). Post-hoc validation score: 53.0.
  \item \textbf{Cluster 2 --- High-density capital hubs}
    (57 zones, 40.1\%): High density ($\log$-density $= 5.60$),
    high sectoral presence ($n=9.70$ sectors), and high capital
    intensity (mean capital $= 15.7$M MAD). Post-hoc validation
    score: 61.3.
\end{itemize}

The monotonic progression of post-hoc validation scores
(35.0 $\rightarrow$ 53.0 $\rightarrow$ 61.3) confirms
economic coherence of the discovered typology
without introducing leakage.

\begin{figure}[htbp]
\centerline{\includegraphics[width=\columnwidth]{artifacts/figures/cluster_profile_heatmap_typology.png}}
\caption{Normalized feature heatmap of the three discovered economic profiles.
  Darker cells indicate higher normalized feature values.}
\label{fig:heatmap}
\end{figure}

\subsection{Supervised Classification Results}

Table~\ref{tab:classifier} reports classifier performance on the held-out
test set. Random~Forest achieved the best macro-F1 (0.955) with
99\% accuracy, selected over LightGBM and XGBoost by cross-validated
macro-F1. These results confirm the generalizability of discovered
profiles to unseen zones.

\begin{table}[htbp]
\caption{Cluster Classifier Performance (Test Set, Typology Mode)}
\label{tab:classifier}
\begin{center}
\begin{tabular}{|l|c|c|}
\hline
\textbf{Model} & \textbf{CV F1-macro} & \textbf{Test F1-macro} \\
\hline
Random Forest & \textbf{0.9587} & \textbf{0.9552} \\
XGBoost       & 0.8961          & ---             \\
LightGBM      & 0.8860          & ---             \\
\hline
\multicolumn{2}{|l|}{Test Accuracy (Random Forest)} & \textbf{0.9655} \\
\hline
\end{tabular}
\end{center}
\end{table}

\subsection{Feature Importance and Explainability}

Fig.~\ref{fig:density} illustrates the spatial distribution of
business density across the study region. Feature importance analysis
identifies the following top predictors of cluster assignment
(descending order): \texttt{active\_rate} (0.201),
\texttt{density\_log} (0.114), \texttt{entity\_count\_real} (0.113),
\texttt{n\_sectors\_present} (0.105), \texttt{entity\_count\_total} (0.100),
\texttt{sarl\_count} (0.093), and \texttt{formal\_ratio} (0.065).
Spatial coordinates (\texttt{cell\_lat}, \texttt{cell\_lon}) contribute
marginally, indicating that the profile structure is driven by
economic intensity rather than geographic location per se.

\begin{figure}[htbp]
\centerline{\includegraphics[width=\columnwidth]{artifacts/figures/morocco_density_heatmap_typology.png}}
\caption{Spatial density heatmap of formal business entities
  across the Casablanca--Settat--El~Jadida region (typology mode).
  Darker areas indicate higher log-density.}
\label{fig:density}
\end{figure}

This explainability pattern is economically coherent:
the classifier separates profiles primarily by \textit{economic intensity}
(density, entity counts), \textit{business vitality} (activity rate),
\textit{structural complexity} (number of sectors), and
\textit{formalization} (SARL count, formal ratio)---all interpretable
dimensions of a spatial economic typology.

% -------------------------------------------------------
\section{Discussion}
\label{sec:discussion}

\subsection{Scientific Validity and Anti-Leakage Design}

A central design choice in this work is the exclusion of the
synthetic attractiveness score from all model inputs. This prevents
the circular reasoning that would arise from predicting a score
built from the same variables used for clustering
\cite{feldmeyer2020using,sakti2024machine}. The score is retained
exclusively for post-hoc validation, playing a role analogous to
an external descriptive signal rather than a ground truth.

\subsection{Practical Implications}

The discovered profiles have direct policy relevance.
Cluster~2 (\textit{High-density capital hubs}) zones can be targeted
for advanced commercial infrastructure and high-value service investment.
Cluster~1 (\textit{Active emerging zones}) represents a strategic
intermediate population---economically active but capital-constrained---
well-suited for support programs targeting SME formalization.
Cluster~0 (\textit{Underdeveloped early-stage ecosystems}) zones indicate
areas where baseline economic stimulation and infrastructure
investment may be needed.

The Random~Forest classifier enables operational deployment:
new zones from other Moroccan regions can be assigned to economic
profiles by computing their feature vector, without requiring
a complete re-clustering of the full dataset.

\subsection{Limitations}

The main limitations of this study are as follows.
The analysis is geographically constrained to the
Casablanca--Settat--El~Jadida region; generalization to other
Moroccan regions requires validation against independent data.
Cluster labels are latent constructs, not official administrative
economic categories, and should be interpreted as descriptive
archetypes rather than absolute classifications.
The unsupervised solution is sensitive to feature scaling and
hyperparameter choices; multi-model selection mitigates but does
not eliminate this sensitivity.
Finally, data are treated as a static snapshot; temporal dynamics
and macroeconomic shocks are not modeled.

% -------------------------------------------------------
\section{Conclusion}
\label{sec:conclusion}

This paper presented a data-driven pipeline for discovering,
characterizing, and assigning economic spatial profiles in Morocco
using unsupervised and supervised machine learning.
Applied to 142 geographic zones in the Casablanca--Settat--El~Jadida
region, Spectral Clustering ($k=3$) revealed three interpretable
economic archetypes that are economically coherent, structurally stable,
and predictively generalizable. A Random~Forest classifier achieved
a test macro-F1 of 0.955, confirming that discovered profiles
can be assigned to new zones with high reliability.
Feature importance analysis identified business activity, density,
sectoral diversification, and formalization as the primary structural
drivers of spatial economic differentiation.

This framework is directly applicable to other Moroccan regions
and to comparable data-scarce emerging-economy contexts.
Future work will extend the pipeline to a national multi-regional
scope, incorporate temporal dynamics through panel data,
and enrich the feature set with complementary sources
(OpenStreetMap amenities, mobility data) for richer
spatial economic characterization.

% -------------------------------------------------------
\section*{Acknowledgment}

The authors thank EMSI Marrakech for institutional support
and access to computing resources.

% -------------------------------------------------------
\begin{thebibliography}{00}

\bibitem{xu2023enhancing}
X. Xu, Z. Chen, and S. Chen,
``Enhancing economic competitiveness analysis through machine learning:
Exploring complex urban features,''
\textit{PLoS ONE}, vol.~18, no.~11, p.~e0293303, 2023.

\bibitem{castellana2024exploring}
F. Castellana et~al.,
``Exploring Apulia's regional tourism attractiveness through the lens of
sustainability: A machine learning approach and counterfactual explainability process,''
\textit{Sustainability}, vol.~16, no.~15, p.~6287, 2024.

\bibitem{khalid2023artificial}
S. Khalid, D. Effina, K. R. Khalid, and M. S. Chaabane,
``The artificial intelligence as a decision-making instrument for modeling
and predicting small cities' attractiveness: Evidence from Morocco,''
\textit{Applied Mathematics \& Information Sciences}, vol.~17, no.~5,
pp.~905--913, 2023.

\bibitem{feldmeyer2020using}
D. Feldmeyer, C. Meisch, H. Sauter, and J. Birkmann,
``Using OpenStreetMap data and machine learning to generate socio-economic indicators,''
\textit{ISPRS International Journal of Geo-Information},
vol.~9, no.~9, p.~498, 2020.

\bibitem{chen2025machine}
S. Chen et~al.,
``Machine learning approaches to predicting urban park attractiveness:
Insights from Shanghai and Tokyo,''
\textit{Urban Forestry \& Urban Greening}, vol.~112, p.~128921, 2025.

\bibitem{huang2022machine}
X. Huang, P. Gong, S. Wang, M. White, and B. Zhang,
``Machine learning modeling of vitality characteristics in historical
preservation zones with multi-source data,''
\textit{Buildings}, vol.~12, no.~11, p.~1978, 2022.

\bibitem{kourtit2021safe}
K. Kourtit, M. M. Marinescu Pele, P. Nijkamp, and D. T. Pele,
``Safe cities in the new urban world: A comparative cluster dynamics
analysis through machine learning,''
\textit{Sustainable Cities and Society}, vol.~66, p.~102665, 2021.

\bibitem{agiropoulos2024exploring}
C. Agiropoulos, J. M. Chen, G. Galanos, and T. Poufinas,
``Exploring regional determinants of tourism success in the eurozone:
An unsupervised machine learning approach,''
\textit{Engineering Proceedings}, vol.~68, no.~1, p.~53, 2024.

\bibitem{sakti2024machine}
A. D. Sakti et~al.,
``Machine learning based urban sprawl assessment using integrated
multi-hazard and environmental-economic impact,''
\textit{Scientific Reports}, vol.~14, no.~1, p.~13385, 2024.

\bibitem{giglio2020machine}
S. Giglio, F. Bertacchini, E. Bilotta, and P. Pantano,
``Machine learning and points of interest: Typical tourist Italian cities,''
\textit{Current Issues in Tourism}, vol.~23, no.~13, pp.~1646--1658, 2020.

\bibitem{bi2024collaborative}
J.-W. Bi, T.-Y. Han, and Y. Yao,
``Collaborative forecasting of tourism demand for multiple tourist attractions
with spatial dependence: A combined deep learning model,''
\textit{Tourism Economics}, vol.~30, no.~2, pp.~361--388, 2024.

\bibitem{abang2022using}
A. Z. Abang Abdurahman et~al.,
``Using machine learning to predict visitors to totally protected areas
in Sarawak, Malaysia,''
\textit{Sustainability}, vol.~14, no.~5, p.~2735, 2022.

\end{thebibliography}

\end{document}